\subsection{Continuity of finite spectral subspaces}\label{subsec:spectral-subspace-continuity}
\paraid{p-f52fe507fe4d}
The distance operators for different compact spaces act on different Hilbert spaces.
We extend the distance kernels to one compact metric space
$\MetricSpace{\AmbientCarrier}{\AmbientMetric}$
to obtain operators on the fixed Banach space
$\ContinuousFunctions(\AmbientCarrier ,\ComplexNumbers)$.
We apply the spectral inclusion and spectral projection convergence theorems
of Anselone and Palmer,
\Cref{thm:spectral-inclusion,thm:riesz-convergence}
in Subsection~\ref{subsec:spectral-preliminaries-perturbation},
to these ambient operators.
By spectral inclusion,
the gaps at the chosen real numbers
$\pm\SpectralCutoff$
persist for all sufficiently large indices.
By convergence of the spectral projections,
the retained spectral subspaces
have the same dimension for all sufficiently large indices.

\paraid{p-040fa54a5428}
By \Cref{lem:parameter-integrals},
the parameter-dependent integrals used below converge uniformly.

\paraid{p-f33149a827cc}
The next proposition compares the generalized eigenspaces of the ambient
kernel operator with those of the self-adjoint distance operator.

\begin{proposition}\label{lem:ambient-eigenspaces}
\paraid{p-c756883ac07c}
Let
$\ComparisonCarrier$
be a nonempty compact subset of a compact metric space
$\MetricSpace{\AmbientCarrier}{\AmbientMetric}$,
and put
\[
 \ComparisonMetric=\AmbientMetric|_{\ComparisonCarrier\times\ComparisonCarrier},
\]
and
\[
 \DiameterBound=\diam\MetricSpace{\AmbientCarrier}{\AmbientMetric}.
\]
Regard
$\DistanceOperator{\ComparisonCarrier}{\ComparisonMetric}$
as its complex-linear extension
\[
 \DistanceOperator{\ComparisonCarrier}{\ComparisonMetric}\colon
 \LebesgueSpace^2(\ComparisonCarrier,\SelectedMeasure{\ComparisonCarrier}{\ComparisonMetric};\ComplexNumbers)
 \to
 \LebesgueSpace^2(\ComparisonCarrier,\SelectedMeasure{\ComparisonCarrier}{\ComparisonMetric};\ComplexNumbers).
\]
Define
$\RestrictionOperator _\ComparisonCarrier \colon\ContinuousFunctions(\AmbientCarrier ,\ComplexNumbers)\to
                       \LebesgueSpace^2(\ComparisonCarrier ,\SelectedMeasure{\ComparisonCarrier}{\ComparisonMetric};\ComplexNumbers)$
as follows.
For every
$\HilbertFunction\in\ContinuousFunctions(\AmbientCarrier,\ComplexNumbers)$,
let
$\RestrictionOperator_\ComparisonCarrier\HilbertFunction$
be the equivalence class of
$\HilbertFunction|_\ComparisonCarrier$.
Define
$\KernelExtensionOperator _\ComparisonCarrier \colon\LebesgueSpace^2(\ComparisonCarrier ,\SelectedMeasure{\ComparisonCarrier}{\ComparisonMetric};\ComplexNumbers)
                       \to\ContinuousFunctions(\AmbientCarrier ,\ComplexNumbers)$
as follows.
For every
$\HilbertTestFunction\in\LebesgueSpace^2(\ComparisonCarrier,\SelectedMeasure{\ComparisonCarrier}{\ComparisonMetric};\ComplexNumbers)$,
using
$\SecondVector\in\ComparisonCarrier$
as the integration variable,
define for every
$\IntegrationPoint\in\AmbientCarrier$
\[
 (\KernelExtensionOperator_\ComparisonCarrier\HilbertTestFunction)(\IntegrationPoint)
 =\int_\ComparisonCarrier
 \AmbientMetric(\IntegrationPoint,\SecondVector)\HilbertTestFunction(\SecondVector)
 \,\IntegrationDifferential\SelectedMeasure{\ComparisonCarrier}{\ComparisonMetric}(\SecondVector).
\]
The integral depends only on the equivalence class of
$\HilbertTestFunction$.
The bounded continuous kernel defines a continuous function on
$\AmbientCarrier$.
Set
\[
 \CarrierAmbientOperator_\ComparisonCarrier
 =\KernelExtensionOperator_\ComparisonCarrier\RestrictionOperator_\ComparisonCarrier
 \colon\ContinuousFunctions(\AmbientCarrier,\ComplexNumbers)
 \to\ContinuousFunctions(\AmbientCarrier,\ComplexNumbers).
\]
Then
$\CarrierAmbientOperator_\ComparisonCarrier$
is compact.
For every
$\Eigenvalue\in\ComplexNumbers\setminus\{0\}$,
the generalized eigenspaces satisfy
\[
 \GeneralizedEigenspace{\CarrierAmbientOperator_\ComparisonCarrier}{\Eigenvalue}
 =\ker(\CarrierAmbientOperator_\ComparisonCarrier-\Eigenvalue\IdentityMap),
\]
and
\[
 \GeneralizedEigenspace{\DistanceOperator{\ComparisonCarrier}{\ComparisonMetric}}{\Eigenvalue}
 =\ker(\DistanceOperator{\ComparisonCarrier}{\ComparisonMetric}-\Eigenvalue\IdentityMap).
\]
Moreover, for every
$\Eigenvalue\in\ComplexNumbers\setminus\{0\}$
and every integer
$\RootOrder\geq1$,
\[
 \ker\bigl((\CarrierAmbientOperator_\ComparisonCarrier-\Eigenvalue\IdentityMap)^\RootOrder\bigr)
 =\ker(\CarrierAmbientOperator_\ComparisonCarrier-\Eigenvalue\IdentityMap),
\]
and
\[
 \ker\bigl((\DistanceOperator{\ComparisonCarrier}{\ComparisonMetric}-\Eigenvalue\IdentityMap)^\RootOrder\bigr)
 =\ker(\DistanceOperator{\ComparisonCarrier}{\ComparisonMetric}-\Eigenvalue\IdentityMap).
\]
Restriction induces the linear isomorphism
\[
 \RestrictionOperator_\ComparisonCarrier\colon
 \GeneralizedEigenspace{\CarrierAmbientOperator_\ComparisonCarrier}{\Eigenvalue}
 \to
 \GeneralizedEigenspace{\DistanceOperator{\ComparisonCarrier}{\ComparisonMetric}}{\Eigenvalue},
\]
Its inverse is the restricted map
\[
 \Eigenvalue^{-1}\KernelExtensionOperator_\ComparisonCarrier\colon
 \GeneralizedEigenspace{\DistanceOperator{\ComparisonCarrier}{\ComparisonMetric}}{\Eigenvalue}
 \to\GeneralizedEigenspace{\CarrierAmbientOperator_\ComparisonCarrier}{\Eigenvalue}.
\]
Moreover,
\[
 \spec(\CarrierAmbientOperator_\ComparisonCarrier)\setminus\{0\}
 =\spec(\DistanceOperator{\ComparisonCarrier}{\ComparisonMetric})\setminus\{0\}
 \subset[-\DiameterBound,\DiameterBound].
\]
\end{proposition}

\begin{proof}
\textbf{Step 1. Compactness and factorization.}
\paraid{p-03d00ab877e8}
Since
$\SelectedMeasure{\ComparisonCarrier}{\ComparisonMetric}$
has total mass one,
for every
$\HilbertFunction\in\ContinuousFunctions(\AmbientCarrier,\ComplexNumbers)$
we have
\[
 \norm{\RestrictionOperator_\ComparisonCarrier\HilbertFunction}_2
 \leq\norm{\HilbertFunction}_\infty.
\]
For every
$\HilbertTestFunction\in\LebesgueSpace^2(\ComparisonCarrier,\SelectedMeasure{\ComparisonCarrier}{\ComparisonMetric};\ComplexNumbers)$
and
$\IntegrationPoint,\IntegrationPoint_0\in\AmbientCarrier$,
the Cauchy--Schwarz inequality yields
\[
 \norm{\KernelExtensionOperator_\ComparisonCarrier\HilbertTestFunction}_\infty
 \leq\DiameterBound\norm{\HilbertTestFunction}_2,
\]
and
\[
 |\KernelExtensionOperator_\ComparisonCarrier\HilbertTestFunction(\IntegrationPoint)
  -\KernelExtensionOperator_\ComparisonCarrier\HilbertTestFunction(\IntegrationPoint_0)|
 \leq\AmbientMetric(\IntegrationPoint,\IntegrationPoint_0)
       \norm{\HilbertTestFunction}_2.
\]
Thus
$\RestrictionOperator_\ComparisonCarrier$
and
$\KernelExtensionOperator_\ComparisonCarrier$
are bounded,
and the Arzel\`a--Ascoli theorem (\Cref{thm:arzela-ascoli}) implies that
$\KernelExtensionOperator_\ComparisonCarrier$
is compact.
Consequently,
$\CarrierAmbientOperator_\ComparisonCarrier$
is compact.
By their definitions,
\begin{equation}\label{eq:operator-factorization-1}
 \CarrierAmbientOperator _\ComparisonCarrier =\KernelExtensionOperator _\ComparisonCarrier \RestrictionOperator _\ComparisonCarrier ,
\end{equation}
and
\begin{equation}\label{eq:operator-factorization-2}
 \DistanceOperator{\ComparisonCarrier}{\ComparisonMetric}=\RestrictionOperator _\ComparisonCarrier \KernelExtensionOperator _\ComparisonCarrier ,
\end{equation}
and
\begin{equation}\label{eq:operator-factorization-3}
 \RestrictionOperator _\ComparisonCarrier  \CarrierAmbientOperator _\ComparisonCarrier =\DistanceOperator{\ComparisonCarrier}{\ComparisonMetric}\RestrictionOperator _\ComparisonCarrier .
\end{equation}

\paraid{p-55a3ce893fd0}
The intertwining identity is displayed in Figure~\ref{fig:operator-comparison}.
\begin{figure}[H]
\centering
\begin{tikzcd}[column sep=large,row sep=large]
 \ContinuousFunctions(\AmbientCarrier,\ComplexNumbers)
 \arrow[r,"\CarrierAmbientOperator_\ComparisonCarrier"]
 \arrow[d,"\RestrictionOperator_\ComparisonCarrier"'] &
 \ContinuousFunctions(\AmbientCarrier,\ComplexNumbers)
 \arrow[d,"\RestrictionOperator_\ComparisonCarrier"] \\
 \LebesgueSpace^2(\ComparisonCarrier,\SelectedMeasure{\ComparisonCarrier}{\ComparisonMetric};\ComplexNumbers)
 \arrow[r,"\DistanceOperator{\ComparisonCarrier}{\ComparisonMetric}"']
 \arrow[ur,"\KernelExtensionOperator_\ComparisonCarrier" description] &
 \LebesgueSpace^2(\ComparisonCarrier,\SelectedMeasure{\ComparisonCarrier}{\ComparisonMetric};\ComplexNumbers)
\end{tikzcd}
\caption{Restriction intertwines the ambient operator
$\CarrierAmbientOperator_\ComparisonCarrier$
and the distance operator
$\DistanceOperator{\ComparisonCarrier}{\ComparisonMetric}$.
The factorizations
$\CarrierAmbientOperator_\ComparisonCarrier=\KernelExtensionOperator_\ComparisonCarrier\RestrictionOperator_\ComparisonCarrier$
and
$\DistanceOperator{\ComparisonCarrier}{\ComparisonMetric}=\RestrictionOperator_\ComparisonCarrier\KernelExtensionOperator_\ComparisonCarrier$
explain the comparison of their nonzero spectral subspaces.}
\label{fig:operator-comparison}
\end{figure}


\textbf{Step 2. Generalized eigenspaces and restriction.}
\paraid{p-3b9fc77c0141}
Fix
$\Eigenvalue\in\ComplexNumbers\setminus\{0\}$
and an integer
$\RootOrder\geq1$.
By the self-adjoint spectral theorem (\Cref{prop:spectral-basics}),
\begin{equation}\label{eq:self-adjoint-power-kernel}
 \ker\bigl((\DistanceOperator{\ComparisonCarrier}{\ComparisonMetric}
               -\Eigenvalue\IdentityMap)^\RootOrder\bigr)
 =\ker(\DistanceOperator{\ComparisonCarrier}{\ComparisonMetric}
               -\Eigenvalue\IdentityMap).
\end{equation}
We prove the equality
\begin{equation}\label{eq:restricted-power-kernel}
 \RestrictionOperator_\ComparisonCarrier
 \left(\ker\bigl((\CarrierAmbientOperator_\ComparisonCarrier
                    -\Eigenvalue\IdentityMap)^\RootOrder\bigr)\right)
 =\ker(\DistanceOperator{\ComparisonCarrier}{\ComparisonMetric}
                    -\Eigenvalue\IdentityMap)
\end{equation}
by both inclusions.
Take
$\HilbertTestFunction\in
 \RestrictionOperator_\ComparisonCarrier
 \left(\ker\bigl((\CarrierAmbientOperator_\ComparisonCarrier
                 -\Eigenvalue\IdentityMap)^\RootOrder\bigr)\right)$.
Then there is
$\HilbertFunction\in\ker\bigl((\CarrierAmbientOperator_\ComparisonCarrier
 -\Eigenvalue\IdentityMap)^\RootOrder\bigr)$
with
$\HilbertTestFunction=\RestrictionOperator_\ComparisonCarrier\HilbertFunction$.
By \eqref{eq:operator-factorization-3},
\[
 (\DistanceOperator{\ComparisonCarrier}{\ComparisonMetric}-\Eigenvalue\IdentityMap)^\RootOrder
 \HilbertTestFunction
 =\RestrictionOperator_\ComparisonCarrier
  (\CarrierAmbientOperator_\ComparisonCarrier-\Eigenvalue\IdentityMap)^\RootOrder
  \HilbertFunction=0.
\]
Equation \eqref{eq:self-adjoint-power-kernel} therefore implies
$\HilbertTestFunction\in\ker(\DistanceOperator{\ComparisonCarrier}{\ComparisonMetric}
 -\Eigenvalue\IdentityMap)$.

\paraid{p-8a2730244eab}
Conversely,
take
$\HilbertTestFunction\in\ker(\DistanceOperator{\ComparisonCarrier}{\ComparisonMetric}
 -\Eigenvalue\IdentityMap)$
and put
$\HilbertFunction=\Eigenvalue^{-1}\KernelExtensionOperator_\ComparisonCarrier\HilbertTestFunction$.
Using \eqref{eq:operator-factorization-1} and \eqref{eq:operator-factorization-2},
we obtain
\begin{equation}\label{eq:eigenfunction-extension-1}
 \RestrictionOperator_\ComparisonCarrier\HilbertFunction
 =\Eigenvalue^{-1}\DistanceOperator{\ComparisonCarrier}{\ComparisonMetric}
    \HilbertTestFunction=\HilbertTestFunction,
\end{equation}
and
\begin{equation}\label{eq:eigenfunction-extension-2}
 \CarrierAmbientOperator_\ComparisonCarrier\HilbertFunction
 =\KernelExtensionOperator_\ComparisonCarrier\HilbertTestFunction
 =\Eigenvalue\HilbertFunction.
\end{equation}
Thus
$\HilbertFunction\in\ker\bigl((\CarrierAmbientOperator_\ComparisonCarrier
 -\Eigenvalue\IdentityMap)^\RootOrder\bigr)$,
and \eqref{eq:eigenfunction-extension-1} implies
$\HilbertTestFunction\in
 \RestrictionOperator_\ComparisonCarrier
 \left(\ker\bigl((\CarrierAmbientOperator_\ComparisonCarrier
                 -\Eigenvalue\IdentityMap)^\RootOrder\bigr)\right)$.
This proves \eqref{eq:restricted-power-kernel}.

\paraid{p-7ddba0188af2}
The restriction operator
$\RestrictionOperator_\ComparisonCarrier$
is injective on
$\ker\bigl((\CarrierAmbientOperator_\ComparisonCarrier-\Eigenvalue\IdentityMap)^\RootOrder\bigr)$.
To see this,
take
$\HilbertFunction$
in this kernel with
$\RestrictionOperator_\ComparisonCarrier\HilbertFunction=0$.
Then \eqref{eq:operator-factorization-1} implies
$\CarrierAmbientOperator_\ComparisonCarrier\HilbertFunction=0$,
so
\[
 0=(\CarrierAmbientOperator_\ComparisonCarrier-\Eigenvalue\IdentityMap)^\RootOrder
     \HilbertFunction=(-\Eigenvalue)^\RootOrder\HilbertFunction.
\]
Since
$\Eigenvalue\ne0$,
we have
$\HilbertFunction=0$.

\paraid{p-60e424cf6d26}
We next prove
\begin{equation}\label{eq:ambient-power-kernel}
 \ker\bigl((\CarrierAmbientOperator_\ComparisonCarrier-\Eigenvalue\IdentityMap)^\RootOrder\bigr)
 =\ker(\CarrierAmbientOperator_\ComparisonCarrier-\Eigenvalue\IdentityMap).
\end{equation}
Take
$\HilbertFunction\in\ker\bigl((\CarrierAmbientOperator_\ComparisonCarrier
 -\Eigenvalue\IdentityMap)^\RootOrder\bigr)$
and set
\[
 \HilbertTestFunction=\RestrictionOperator_\ComparisonCarrier\HilbertFunction,
\]
and
\[
 \LiftedEigenfunction=\Eigenvalue^{-1}\KernelExtensionOperator_\ComparisonCarrier\HilbertTestFunction.
\]
Equations \eqref{eq:restricted-power-kernel},
\eqref{eq:eigenfunction-extension-1} and \eqref{eq:eigenfunction-extension-2}
imply
\[
 \LiftedEigenfunction\in\ker(\CarrierAmbientOperator_\ComparisonCarrier-\Eigenvalue\IdentityMap),
\]
and
\[
 \RestrictionOperator_\ComparisonCarrier\LiftedEigenfunction
 =\RestrictionOperator_\ComparisonCarrier\HilbertFunction.
\]
Both
$\HilbertFunction$
and
$\LiftedEigenfunction$
belong to
\[
 \ker\bigl((\CarrierAmbientOperator_\ComparisonCarrier
 -\Eigenvalue\IdentityMap)^\RootOrder\bigr),
\]
the kernel on the left of \eqref{eq:ambient-power-kernel}.
Injectivity of restriction on this kernel implies
$\HilbertFunction=\LiftedEigenfunction$,
so
$\HilbertFunction\in\ker(\CarrierAmbientOperator_\ComparisonCarrier-\Eigenvalue\IdentityMap)$.
Conversely,
if
$\HilbertFunction\in\ker(\CarrierAmbientOperator_\ComparisonCarrier-\Eigenvalue\IdentityMap)$,
then
\[
 (\CarrierAmbientOperator_\ComparisonCarrier-\Eigenvalue\IdentityMap)^\RootOrder\HilbertFunction
 =(\CarrierAmbientOperator_\ComparisonCarrier-\Eigenvalue\IdentityMap)^{\RootOrder-1}0=0.
\]
Hence
$\HilbertFunction\in\ker\bigl((\CarrierAmbientOperator_\ComparisonCarrier
 -\Eigenvalue\IdentityMap)^\RootOrder\bigr)$.
This proves \eqref{eq:ambient-power-kernel}.
Taking the union over
$\RootOrder\geq1$
in \eqref{eq:self-adjoint-power-kernel} and \eqref{eq:ambient-power-kernel}
identifies the corresponding generalized eigenspaces.
In particular,
restriction induces the linear isomorphism
\[
 \RestrictionOperator_\ComparisonCarrier\colon
 \ker(\CarrierAmbientOperator_\ComparisonCarrier-\Eigenvalue\IdentityMap)
 \to\ker(\DistanceOperator{\ComparisonCarrier}{\ComparisonMetric}-\Eigenvalue\IdentityMap),
\]
Its inverse is the restricted map
\[
 \Eigenvalue^{-1}
 \left.\KernelExtensionOperator_\ComparisonCarrier\right|
 _{\ker(\DistanceOperator{\ComparisonCarrier}{\ComparisonMetric}-\Eigenvalue\IdentityMap)}\colon
 \ker(\DistanceOperator{\ComparisonCarrier}{\ComparisonMetric}-\Eigenvalue\IdentityMap)
 \to\ker(\CarrierAmbientOperator_\ComparisonCarrier-\Eigenvalue\IdentityMap).
\]

\textbf{Step 3. Equality of the nonzero spectra.}
\paraid{p-7bd77f5fa755}
Both operators are compact,
so the Riesz--Schauder theorem (\Cref{thm:compact-spectrum}) ensures
that every nonzero spectral value is an eigenvalue.
By \Cref{lem:distance-operator},
$\norm{\DistanceOperator{\ComparisonCarrier}{\ComparisonMetric}}
 \leq\diam\MetricSpace{\ComparisonCarrier}{\ComparisonMetric}
 \leq\DiameterBound$.
The standard spectral bound for bounded operators
\cite[Theorem~14.3]{Muscat2024}
implies
\[
 \spec(\DistanceOperator{\ComparisonCarrier}{\ComparisonMetric})
 \subset
 \{\Eigenvalue\in\ComplexNumbers\mid
   |\Eigenvalue|\leq
   \norm{\DistanceOperator{\ComparisonCarrier}{\ComparisonMetric}}\}.
\]
Since
$\DistanceOperator{\ComparisonCarrier}{\ComparisonMetric}$
is self-adjoint, its spectrum is real. Thus
\[
 \spec(\CarrierAmbientOperator _\ComparisonCarrier )\setminus\{0\}
 =\spec(\DistanceOperator{\ComparisonCarrier}{\ComparisonMetric})\setminus\{0\}
 \subset[-\DiameterBound,\DiameterBound].
\]
This proves the equality and the bound for the nonzero spectra,
and finishes the proof.
\end{proof}

\paraid{p-6a4c7f8bfb4d}
We first compare the ambient Riesz projection with the orthogonal projection
for a fixed compact carrier.
We then prove convergence of the ambient projections as the carrier varies.

\begin{proposition}\label{lem:projection-restriction}
\paraid{p-7625e731f8c2}
Let
$\ComparisonCarrier$
be a nonempty compact subset of a compact metric space
$\MetricSpace{\AmbientCarrier}{\AmbientMetric}$,
and put
$\ComparisonMetric=\AmbientMetric|_{\ComparisonCarrier\times\ComparisonCarrier}$.
Use the operators of \Cref{lem:ambient-eigenspaces}.
Fix
$\SpectralCutoff>0$
with
$\pm\SpectralCutoff\notin\spec(\DistanceOperator{\ComparisonCarrier}{\ComparisonMetric})$.
Let
$\AmbientProjection_\ComparisonCarrier$
be the Riesz projection of
$\CarrierAmbientOperator_\ComparisonCarrier$
associated with the eigenvalues of absolute value greater than
$\SpectralCutoff$,
as in Definition~\ref{def:riesz-projection}.
Then this projection preserves real functions,
and restriction induces the complex-linear isomorphism
\[
 \RestrictionOperator_\ComparisonCarrier\colon
 \range\AmbientProjection_\ComparisonCarrier
 \to\SpectralSubspace{\ComparisonCarrier}{\SpectralCutoff}
   +\ImaginaryUnit\SpectralSubspace{\ComparisonCarrier}{\SpectralCutoff}.
\]
Use
$\SpectralProjection{\ComparisonCarrier}{\SpectralCutoff}$
also for the complex-linear extension of the orthogonal projection.
For every
$\HilbertFunction\in\ContinuousFunctions(\AmbientCarrier,\ComplexNumbers)$
we have
\[
 (\AmbientProjection_\ComparisonCarrier\HilbertFunction)|_\ComparisonCarrier
 =\SpectralProjection{\ComparisonCarrier}{\SpectralCutoff}
   (\HilbertFunction|_\ComparisonCarrier).
\]
\end{proposition}

\begin{proof}
\paraid{p-8812db27373e}
By \Cref{lem:ambient-eigenspaces},
\begin{equation}\label{eq:fixed-carrier-nonzero-spectra}
 \spec(\CarrierAmbientOperator_\ComparisonCarrier)\setminus\{0\}
 =\spec(\DistanceOperator{\ComparisonCarrier}{\ComparisonMetric})\setminus\{0\}
 \subset\RealNumbers.
\end{equation}
Choose a bounded open set
$\SpectralRegion\subset\ComplexNumbers$
that is a union of two rectangles symmetric about the real axis and satisfies
$0\notin\overline{\SpectralRegion}$.
By \eqref{eq:fixed-carrier-nonzero-spectra}
and the hypothesis
$\pm\SpectralCutoff\notin\spec(\DistanceOperator{\ComparisonCarrier}{\ComparisonMetric})$,
we can choose
$\SpectralRegion$
so that for each operator
$\CompactOperator\in\{\CarrierAmbientOperator_\ComparisonCarrier,
\DistanceOperator{\ComparisonCarrier}{\ComparisonMetric}\}$,
\begin{equation}\label{eq:fixed-carrier-selected-spectrum}
 \SpectralRegion\cap\spec(\CompactOperator)
 =\{\Eigenvalue\in\spec(\CompactOperator)\mid|\Eigenvalue|>\SpectralCutoff\},
\end{equation}
and
\[
 \partial\SpectralRegion\cap\spec(\CompactOperator)=\emptyset.
\]
Orient
$\SpectralContour=\partial\SpectralRegion$
positively.
By Definition~\ref{def:riesz-projection}
and \Cref{lem:riesz-projection-properties},
we have
\[
 \AmbientProjection_\ComparisonCarrier
 =\RieszProjection{\CarrierAmbientOperator_\ComparisonCarrier}{\SpectralContour}.
\]
By \Cref{lem:riesz-range-generalized-eigenspace},
\begin{equation}\label{eq:fixed-carrier-riesz-range}
 \range\AmbientProjection_\ComparisonCarrier
 =\bigoplus_{\Eigenvalue\in\SpectralRegion\cap\spec(\CarrierAmbientOperator_\ComparisonCarrier)}
 \GeneralizedEigenspace{\CarrierAmbientOperator_\ComparisonCarrier}{\Eigenvalue}.
\end{equation}
For each
$\Eigenvalue\in\SpectralRegion\cap\spec(\CarrierAmbientOperator_\ComparisonCarrier)$,
\Cref{lem:ambient-eigenspaces} shows that the restriction operator
$\RestrictionOperator_\ComparisonCarrier$
maps
$\GeneralizedEigenspace{\CarrierAmbientOperator_\ComparisonCarrier}{\Eigenvalue}$
isomorphically onto
$\ker(\DistanceOperator{\ComparisonCarrier}{\ComparisonMetric}-\Eigenvalue\IdentityMap)$.
Taking the direct sum in \eqref{eq:fixed-carrier-riesz-range}
and using \eqref{eq:fixed-carrier-selected-spectrum},
we obtain the isomorphism
\[
 \RestrictionOperator_\ComparisonCarrier\colon
 \range\AmbientProjection_\ComparisonCarrier
 \to\bigoplus_{\substack{\Eigenvalue\in\spec(\DistanceOperator{\ComparisonCarrier}{\ComparisonMetric})\\
 |\Eigenvalue|>\SpectralCutoff}}
 \ker(\DistanceOperator{\ComparisonCarrier}{\ComparisonMetric}-\Eigenvalue\IdentityMap).
\]
By the definition of
$\SpectralSubspace{\ComparisonCarrier}{\SpectralCutoff}$,
this direct sum is
$\SpectralSubspace{\ComparisonCarrier}{\SpectralCutoff}
+\ImaginaryUnit\SpectralSubspace{\ComparisonCarrier}{\SpectralCutoff}$.
We prove the identity relating projection and restriction for
an arbitrary continuous function.
For every
$\ResolventParameter\in\SpectralContour$,
\eqref{eq:operator-factorization-3} implies
\[
 (\ResolventParameter\IdentityMap-\DistanceOperator{\ComparisonCarrier}{\ComparisonMetric})
 \RestrictionOperator_\ComparisonCarrier
 =\RestrictionOperator_\ComparisonCarrier
 (\ResolventParameter\IdentityMap-\CarrierAmbientOperator_\ComparisonCarrier).
\]
For
$\ResolventParameter\in\SpectralContour$,
both operators in parentheses are invertible.
Multiplying by their inverses on the left and on the right,
respectively,
we obtain
\begin{equation}\label{eq:resolvent-restriction}
 \RestrictionOperator_\ComparisonCarrier
 (\ResolventParameter\IdentityMap-\CarrierAmbientOperator_\ComparisonCarrier)^{-1}
 =(\ResolventParameter\IdentityMap-\DistanceOperator{\ComparisonCarrier}{\ComparisonMetric})^{-1}
 \RestrictionOperator_\ComparisonCarrier.
\end{equation}
Both sides are bounded maps from
$\ContinuousFunctions(\AmbientCarrier,\ComplexNumbers)$
to
$\LebesgueSpace^2(\ComparisonCarrier,\SelectedMeasure{\ComparisonCarrier}{\ComparisonMetric};\ComplexNumbers)$.
Both resolvents in \eqref{eq:resolvent-restriction} depend continuously on
$\ResolventParameter\in\SpectralContour$
in the operator norm.
The line integrals can therefore be computed as limits of Riemann sums.
Since
$\RestrictionOperator_\ComparisonCarrier$
is bounded and linear,
it commutes with these sums and their limits.
For every
$\HilbertFunction\in\ContinuousFunctions(\AmbientCarrier,\ComplexNumbers)$,
we consequently obtain
\begin{equation}\label{eq:restricted-contour-integral}
 \begin{aligned}
 \RestrictionOperator_\ComparisonCarrier\AmbientProjection_\ComparisonCarrier\HilbertFunction
 &=\RestrictionOperator_\ComparisonCarrier
   \left(\frac{1}{2\CircleConstant\ImaginaryUnit}
   \int_\SpectralContour
   (\ResolventParameter\IdentityMap-\CarrierAmbientOperator_\ComparisonCarrier)^{-1}
   \HilbertFunction\,\IntegrationDifferential\ResolventParameter\right)\\
 &=\frac{1}{2\CircleConstant\ImaginaryUnit}
   \int_\SpectralContour\RestrictionOperator_\ComparisonCarrier
   (\ResolventParameter\IdentityMap-\CarrierAmbientOperator_\ComparisonCarrier)^{-1}
   \HilbertFunction\,\IntegrationDifferential\ResolventParameter\\
 &=\frac{1}{2\CircleConstant\ImaginaryUnit}
   \int_\SpectralContour
   (\ResolventParameter\IdentityMap-\DistanceOperator{\ComparisonCarrier}{\ComparisonMetric})^{-1}
   \RestrictionOperator_\ComparisonCarrier\HilbertFunction
   \,\IntegrationDifferential\ResolventParameter.
 \end{aligned}
\end{equation}
In \eqref{eq:restricted-contour-integral},
the first equality is the definition of
$\AmbientProjection_\ComparisonCarrier$,
the second uses the bounded linearity of
$\RestrictionOperator_\ComparisonCarrier$,
and the third uses \eqref{eq:resolvent-restriction}.

\paraid{p-b3df5075d3de}
We next prove that the last line integral in
\eqref{eq:restricted-contour-integral} equals
$\SpectralProjection{\ComparisonCarrier}{\SpectralCutoff}
 (\RestrictionOperator_\ComparisonCarrier\HilbertFunction)$.
Here
$\SpectralProjection{\ComparisonCarrier}{\SpectralCutoff}$
also denotes the complex-linear extension of the orthogonal projection
defined in \eqref{eq:spectral-projection-definition}.
Let
$\HilbertTestFunction$
be an eigenvector of
$\DistanceOperator{\ComparisonCarrier}{\ComparisonMetric}$
with eigenvalue
$\Eigenvalue$.
For every
$\ResolventParameter\in\SpectralContour$,
\begin{equation}\label{eq:eigenvector-resolvent}
 (\ResolventParameter\IdentityMap-\DistanceOperator{\ComparisonCarrier}{\ComparisonMetric})^{-1}
 \HilbertTestFunction
 =\frac{1}{\ResolventParameter-\Eigenvalue}\HilbertTestFunction.
\end{equation}
Substituting \eqref{eq:eigenvector-resolvent} into the line integral
and using the winding numbers of
$\SpectralContour$,
we obtain
\begin{equation}\label{eq:eigenvector-contour-integral}
 \begin{aligned}
 &\frac{1}{2\CircleConstant\ImaginaryUnit}
 \int_\SpectralContour
 (\ResolventParameter\IdentityMap-\DistanceOperator{\ComparisonCarrier}{\ComparisonMetric})^{-1}
 \HilbertTestFunction\,\IntegrationDifferential\ResolventParameter\\
 &\qquad=
 \left(\frac{1}{2\CircleConstant\ImaginaryUnit}
 \int_\SpectralContour
 \frac{\IntegrationDifferential\ResolventParameter}{\ResolventParameter-\Eigenvalue}\right)
 \HilbertTestFunction
 =
 \begin{cases}
 \HilbertTestFunction,&|\Eigenvalue|>\SpectralCutoff,\\
 0,&|\Eigenvalue|<\SpectralCutoff.
 \end{cases}
 \end{aligned}
\end{equation}
The scalar line integral equals one when
$\Eigenvalue$
is inside
$\SpectralContour$
and zero when it is outside.
The second case includes
$\Eigenvalue=0$,
so the line integral vanishes on
$\ker\DistanceOperator{\ComparisonCarrier}{\ComparisonMetric}$.
Take the closure in the
$\LebesgueSpace^2$-norm.
By the self-adjoint spectral theorem (\Cref{prop:spectral-basics}),
\begin{equation}\label{eq:hilbert-eigenspace-density}
 \begin{aligned}
 &\overline{
   \ker\DistanceOperator{\ComparisonCarrier}{\ComparisonMetric}
   \mathbin{\oplus}
   \bigoplus_{\substack{
     \Eigenvalue\in\spec(\DistanceOperator{\ComparisonCarrier}{\ComparisonMetric})\\
     \Eigenvalue\ne0}}
   \ker(\DistanceOperator{\ComparisonCarrier}{\ComparisonMetric}-\Eigenvalue\IdentityMap)
   }\\
 &\qquad=
 \LebesgueSpace^2(\ComparisonCarrier,\SelectedMeasure{\ComparisonCarrier}{\ComparisonMetric};\ComplexNumbers).
 \end{aligned}
\end{equation}
By \eqref{eq:eigenvector-contour-integral},
\eqref{eq:spectral-projection-definition},
and orthogonality of eigenspaces for distinct eigenvalues,
the operator defined by the line integral and
$\SpectralProjection{\ComparisonCarrier}{\SpectralCutoff}$
agree on each summand in \eqref{eq:hilbert-eigenspace-density}.
By linearity and boundedness,
the density in \eqref{eq:hilbert-eigenspace-density} proves
\begin{equation}\label{eq:hilbert-riesz-projection}
 \frac{1}{2\CircleConstant\ImaginaryUnit}
 \int_\SpectralContour
 (\ResolventParameter\IdentityMap-\DistanceOperator{\ComparisonCarrier}{\ComparisonMetric})^{-1}
 \,\IntegrationDifferential\ResolventParameter
 =\SpectralProjection{\ComparisonCarrier}{\SpectralCutoff}
\end{equation}
as operators on
$\LebesgueSpace^2(\ComparisonCarrier,\SelectedMeasure{\ComparisonCarrier}{\ComparisonMetric};\ComplexNumbers)$.
Substituting \eqref{eq:hilbert-riesz-projection} into
\eqref{eq:restricted-contour-integral},
we obtain
\[
 \RestrictionOperator_\ComparisonCarrier\AmbientProjection_\ComparisonCarrier\HilbertFunction
 =\SpectralProjection{\ComparisonCarrier}{\SpectralCutoff}
    \RestrictionOperator_\ComparisonCarrier\HilbertFunction.
\]
Since
$\RestrictionOperator_\ComparisonCarrier\HilbertFunction=\HilbertFunction|_\ComparisonCarrier$,
we conclude that
\begin{equation}\label{eq:projection-restriction}
 (\AmbientProjection_\ComparisonCarrier\HilbertFunction)|_\ComparisonCarrier
 =\SpectralProjection{\ComparisonCarrier}{\SpectralCutoff}
   (\HilbertFunction|_\ComparisonCarrier).
\end{equation}
On the left we project in
$\ContinuousFunctions(\AmbientCarrier,\ComplexNumbers)$
and then restrict to
$\ComparisonCarrier$.
On the right we first restrict and then apply the projection on
$\LebesgueSpace^2(\ComparisonCarrier,\SelectedMeasure{\ComparisonCarrier}{\ComparisonMetric};\ComplexNumbers)$.
Both sides of \eqref{eq:projection-restriction} are continuous on
$\ComparisonCarrier$.
Since
$\SelectedMeasure{\ComparisonCarrier}{\ComparisonMetric}$
has full support,
their equality in
$\LebesgueSpace^2$
also implies pointwise equality on
$\ComparisonCarrier$.

\paraid{p-4e9939cad1e5}
Finally,
the distance kernels are real and
$\SpectralRegion$
is invariant under complex conjugation.
Conjugating the line integrals and reversing the orientation of
$\SpectralContour$
therefore yields
\[
 \overline{\AmbientProjection_\ComparisonCarrier\HilbertFunction}
 =\AmbientProjection_\ComparisonCarrier\overline{\HilbertFunction}.
\]
Thus
$\AmbientProjection_\ComparisonCarrier$
preserves real functions.
\end{proof}

\begin{proposition}\label{lem:spectral-projections}
\paraid{p-b10cee5c3ca1}
Let
$\BaseCarrier _\SequenceIndex ,\BaseCarrier $
be nonempty compact subsets of a compact metric space
$\MetricSpace{\AmbientCarrier }{\AmbientMetric }$,
with restricted metrics
\[
 \MetricSymbol_\SequenceIndex=\AmbientMetric|_{\BaseCarrier_\SequenceIndex\times\BaseCarrier_\SequenceIndex},
\]
and
\[
 \MetricSymbol=\AmbientMetric|_{\BaseCarrier\times\BaseCarrier},
\]
such that
\[
 \HausdorffDistance{\AmbientMetric }(\BaseCarrier _\SequenceIndex ,\BaseCarrier )\to0,
\]
and let $\iota_\SequenceIndex\colon\BaseCarrier_\SequenceIndex\hookrightarrow\AmbientCarrier$
and $\iota\colon\BaseCarrier\hookrightarrow\AmbientCarrier$ be the inclusions.
Define probabilities on $\AmbientCarrier$ by
$\bar\mu_\SequenceIndex=(\iota_\SequenceIndex)_*\SelectedMeasure{\BaseCarrier_\SequenceIndex}{\MetricSymbol_\SequenceIndex}$
and $\bar\mu=\iota_*\SelectedMeasure{\BaseCarrier}{\MetricSymbol}$.
Fix
$\SpectralCutoff >0$
with
$\SpectralCutoff ,-\SpectralCutoff \notin\spec(\DistanceOperator{\BaseCarrier}{\MetricSymbol})$.
Use the operators of \Cref{lem:ambient-eigenspaces} and set
\[
 \AmbientOperator_\SequenceIndex=\CarrierAmbientOperator_{\BaseCarrier_\SequenceIndex},
\]
and
\[
 \AmbientOperator=\CarrierAmbientOperator_\BaseCarrier.
\]
Then for all sufficiently large
$\SequenceIndex$,
\[
 \pm\SpectralCutoff
 \notin\spec(\DistanceOperator{\BaseCarrier_\SequenceIndex}{\MetricSymbol_\SequenceIndex}),
\]
and
\[
 \dim\SpectralSubspace{\BaseCarrier_\SequenceIndex}{\SpectralCutoff}
 =\dim\SpectralSubspace{\BaseCarrier}{\SpectralCutoff}
 =\SpectralRank.
\]
Moreover,
there are finite-rank Riesz projections
\[
 \AmbientProjection_\SequenceIndex,\AmbientProjection
 \colon\ContinuousFunctions(\AmbientCarrier,\ComplexNumbers)
 \to\ContinuousFunctions(\AmbientCarrier,\ComplexNumbers)
\]
associated with the eigenvalues of absolute value greater than
$\SpectralCutoff$
of
$\AmbientOperator_\SequenceIndex$
and
$\AmbientOperator$,
respectively.
Their ranges are the finite-dimensional sums
\[
 \range\AmbientProjection
 =\bigoplus_{\substack{\Eigenvalue\in\spec(\AmbientOperator)\\|\Eigenvalue|>\SpectralCutoff}}
 \GeneralizedEigenspace{\AmbientOperator}{\Eigenvalue},
\]
and, for every sufficiently large $\SequenceIndex$,
\[
 \range\AmbientProjection_\SequenceIndex
 =\bigoplus_{\substack{\Eigenvalue\in\spec(\AmbientOperator_\SequenceIndex)\\|\Eigenvalue|>\SpectralCutoff}}
 \GeneralizedEigenspace{\AmbientOperator_\SequenceIndex}{\Eigenvalue}.
\]
Both projections preserve real functions and satisfy,
for every
$\HilbertFunction\in\ContinuousFunctions(\AmbientCarrier,\ComplexNumbers)$,
\[
 \norm{\AmbientProjection_\SequenceIndex\HilbertFunction
       -\AmbientProjection\HilbertFunction}_\infty\to0.
\]
\end{proposition}

\begin{proof}
\textbf{Step 1. Strong convergence and collective compactness.}
\paraid{p-424dfa7891fa}
By \Cref{thm:invariant-measures}\ref{item:invariant-measures-continuity} applied to the inclusions
$\iota_\SequenceIndex$
and
$\iota$,
we have
$\bar\mu_\SequenceIndex\to\bar\mu$
weakly.
For every
$\SequenceIndex\in\NonnegativeIntegers$,
$\HilbertFunction\in\ContinuousFunctions(\AmbientCarrier,\ComplexNumbers)$,
and
$\IntegrationPoint\in\AmbientCarrier$,
the definitions imply
\[
\AmbientOperator _\SequenceIndex \HilbertFunction (\IntegrationPoint )=\int_\AmbientCarrier  \AmbientMetric (\IntegrationPoint ,\SecondVector )\HilbertFunction (\SecondVector )\,\IntegrationDifferential \bar\mu_\SequenceIndex(\SecondVector ),
\]
and
\[
\AmbientOperator \HilbertFunction (\IntegrationPoint )=\int_\AmbientCarrier  \AmbientMetric (\IntegrationPoint ,\SecondVector )\HilbertFunction (\SecondVector )\,\IntegrationDifferential \bar\mu(\SecondVector ).
\]
For each fixed
$\HilbertFunction $,
the family
$\{\AmbientMetric (\IntegrationPoint ,\cdot)\HilbertFunction (\cdot)\mid \IntegrationPoint \in \AmbientCarrier \}$
is compact in
$\ContinuousFunctions(\AmbientCarrier ,\ComplexNumbers)$.
By \Cref{lem:parameter-integrals} with parameter
$\IntegrationPoint\in\AmbientCarrier$,
\begin{equation}\label{eq:ambient-strong-convergence}
 \norm{\AmbientOperator _\SequenceIndex \HilbertFunction -\AmbientOperator \HilbertFunction }_\infty\to0.
\end{equation}
Thus
$\AmbientOperator_\SequenceIndex\to\AmbientOperator$
strongly.
Writing
$\DiameterBound =\diam\MetricSpace{\AmbientCarrier }{\AmbientMetric }$,
for every
$\SequenceIndex\in\NonnegativeIntegers$,
$\HilbertFunction\in\ContinuousFunctions(\AmbientCarrier,\ComplexNumbers)$
with
$\norm{\HilbertFunction}_\infty\leq1$,
and
$\IntegrationPoint,\IntegrationPoint_0\in\AmbientCarrier$,
we obtain
\begin{equation}\label{eq:ambient-operator-bounds-1}
 \norm{\AmbientOperator _\SequenceIndex \HilbertFunction }_\infty\leq \DiameterBound ,
\end{equation}
and
\begin{equation}\label{eq:ambient-operator-bounds-2}
 |\AmbientOperator _\SequenceIndex \HilbertFunction (\IntegrationPoint )-\AmbientOperator _\SequenceIndex \HilbertFunction (\IntegrationPoint _0)|\leq \AmbientMetric (\IntegrationPoint ,\IntegrationPoint _0).
\end{equation}
Hence
$\bigcup_\SequenceIndex  \AmbientOperator _\SequenceIndex \{\HilbertFunction \mid\norm{\HilbertFunction }_\infty\leq1\}$
has compact closure by the Arzel\`a--Ascoli theorem (\Cref{thm:arzela-ascoli}).
Thus
$\{\AmbientOperator_\SequenceIndex\}_{\SequenceIndex\in\NonnegativeIntegers}$
is collectively compact.
Since
$\AmbientOperator$
is compact by \Cref{lem:ambient-eigenspaces},
the family
$\{\AmbientOperator_\SequenceIndex-\AmbientOperator\}_{\SequenceIndex\in\NonnegativeIntegers}$
is also collectively compact.

\textbf{Step 2. A common boundary and spectral projections.}
\paraid{p-8737a7c3af89}
By \Cref{lem:ambient-eigenspaces} and the hypothesis on
$\SpectralCutoff$,
we have
\[
 \spec(\AmbientOperator)
 \subset\ComplexNumbers\setminus\{-\SpectralCutoff,\SpectralCutoff\}.
\]
The set
$\ComplexNumbers\setminus\{-\SpectralCutoff,\SpectralCutoff\}$
is open.
The strong convergence in \eqref{eq:ambient-strong-convergence}
and the collective compactness of
$\{\AmbientOperator_\SequenceIndex-\AmbientOperator\}_{\SequenceIndex\in\NonnegativeIntegers}$
verify the hypotheses on operators in the statement of
\Cref{thm:spectral-inclusion}.
Applying that theorem with this open set
$\ComplexNumbers\setminus\{-\SpectralCutoff,\SpectralCutoff\}$,
we obtain an index
$\SequenceIndex_0\in\NonnegativeIntegers$
such that for every
$\SequenceIndex\geq\SequenceIndex_0$
we have
\begin{equation}\label{eq:ambient-cutoff-exclusion}
 \{-\SpectralCutoff,\SpectralCutoff\}
 \cap\bigl(\spec(\AmbientOperator)\cup\spec(\AmbientOperator_\SequenceIndex)\bigr)
 =\emptyset.
\end{equation}
\Cref{lem:ambient-eigenspaces} then implies
\[
 \pm\SpectralCutoff
 \notin\spec(\DistanceOperator{\BaseCarrier_\SequenceIndex}{\MetricSymbol_\SequenceIndex}).
\]
This proves the claimed exclusion of
$\pm\SpectralCutoff$.

\paraid{p-4c4bdee6941a}
We next choose a bounded open set whose boundary encloses the eigenvalues of absolute
value greater than
$\SpectralCutoff$
for both
$\AmbientOperator$
and
$\AmbientOperator_\SequenceIndex$.
Choose
$\SpectralOuterBound >\max\{\DiameterBound ,\SpectralCutoff \}+1$
and let
$\SpectralRegion $
be the union of the two rectangles
\[
 \begin{gathered}
 \{\ResolventParameter \in\ComplexNumbers\mid \SpectralCutoff <\RealPart \ResolventParameter <\SpectralOuterBound ,\ |\ImaginaryPart \ResolventParameter |<1\},\\
 \{\ResolventParameter \in\ComplexNumbers\mid-\SpectralOuterBound <\RealPart \ResolventParameter <-\SpectralCutoff ,\ |\ImaginaryPart \ResolventParameter |<1\}.
 \end{gathered}
\]
Write
$\SpectralContour =\partial\SpectralRegion $
with positive orientation.
For every
$\SequenceIndex\in\NonnegativeIntegers$,
\Cref{lem:ambient-eigenspaces} implies
\begin{equation}\label{eq:ambient-spectral-bound}
 \spec(\AmbientOperator)\cup\spec(\AmbientOperator_\SequenceIndex)
 \subset[-\DiameterBound,\DiameterBound].
\end{equation}
The choice of
$\SpectralOuterBound$
implies
\begin{equation}\label{eq:contour-real-intersection}
 \SpectralContour\cap[-\DiameterBound,\DiameterBound]
 \subset\{-\SpectralCutoff,\SpectralCutoff\}.
\end{equation}
Combining \eqref{eq:contour-real-intersection} with \eqref{eq:ambient-cutoff-exclusion}
and \eqref{eq:ambient-spectral-bound},
we obtain for every
$\SequenceIndex\geq\SequenceIndex_0$
\begin{equation}\label{eq:common-spectral-contour}
 \SpectralContour\cap
 \bigl(\spec(\AmbientOperator)\cup\spec(\AmbientOperator_\SequenceIndex)\bigr)
 =\emptyset.
\end{equation}
Moreover,
\begin{equation}\label{eq:spectral-region-selection-1}
 \spec(\AmbientOperator)\cap\SpectralRegion
 =\{\Eigenvalue\in\spec(\AmbientOperator)\mid|\Eigenvalue|>\SpectralCutoff\},
\end{equation}
and
\begin{equation}\label{eq:spectral-region-selection-2}
 \spec(\AmbientOperator_\SequenceIndex)\cap\SpectralRegion
 =\{\Eigenvalue\in\spec(\AmbientOperator_\SequenceIndex)\mid|\Eigenvalue|>\SpectralCutoff\},
\end{equation}
and
\begin{equation}\label{eq:spectral-region-selection-3}
 0\notin\overline\SpectralRegion.
\end{equation}
By \eqref{eq:common-spectral-contour},
we use the same boundary for both operators
for all sufficiently large indices.
Equations \eqref{eq:spectral-region-selection-1} and \eqref{eq:spectral-region-selection-2}
show that this boundary encloses precisely the eigenvalues with absolute value
greater than
$\SpectralCutoff$.
Figure~\ref{fig:spectral-contour} shows the two components of this boundary.
\begin{figure}[H]
\centering
\begin{tikzpicture}[>=Stealth,x=1cm,y=1cm,font=\small]
\fill[black!6] (-4,-.8) rectangle (-1,.8);
\fill[black!6] (1,-.8) rectangle (4,.8);
\draw[->] (-4.5,0)--(4.7,0) node[right] {$\RealPart\ResolventParameter$};
\draw[->] (0,-1.05)--(0,1.3) node[above] {$\ImaginaryPart\ResolventParameter$};
\draw[thick] (-4,-.8) rectangle (-1,.8);
\draw[thick] (1,-.8) rectangle (4,.8);
\draw[->,thick] (-3.2,-.8)--(-2.5,-.8);
\draw[->,thick] (2,-.8)--(2.7,-.8);
\foreach \x in {-3.5,-2.8,-1.7,1.6,2.5,3.4} \fill (\x,0) circle (1.5pt);
\foreach \x in {-.55,-.3,.25,.5} \fill[black!45] (\x,0) circle (1.2pt);
\node[below left] at (0,0) {$0$};
\node[below] at (-4,-.85) {$-\SpectralOuterBound$};
\node[below] at (-1,-.85) {$-\SpectralCutoff$};
\node[below] at (1,-.85) {$\SpectralCutoff$};
\node[below] at (4,-.85) {$\SpectralOuterBound$};
\node[above] at (-2.5,.85) {$\SpectralContour$};
\node[above] at (2.5,.85) {$\SpectralContour$};
\end{tikzpicture}
\caption{The positively oriented boundary
$\SpectralContour=\partial\SpectralRegion$
encloses the retained positive and negative spectral values.
The dots are schematic.
The gaps at
$\pm\SpectralCutoff$
allow the same boundary to be used for
$\AmbientOperator_\SequenceIndex$
for all sufficiently large
$\SequenceIndex$.}
\label{fig:spectral-contour}
\end{figure}


\paraid{p-341c679243de}
Using Definition~\ref{def:riesz-projection},
we define the projections
$\AmbientProjection,\AmbientProjection_\SequenceIndex
 \colon\ContinuousFunctions(\AmbientCarrier,\ComplexNumbers)
 \to\ContinuousFunctions(\AmbientCarrier,\ComplexNumbers)$
by
\[
 \AmbientProjection=\RieszProjection{\AmbientOperator}{\SpectralContour},
\]
and
\[
 \AmbientProjection_\SequenceIndex
 =\RieszProjection{\AmbientOperator_\SequenceIndex}{\SpectralContour}.
\]
For
$\HilbertFunction\in\ContinuousFunctions(\AmbientCarrier,\ComplexNumbers)$
and each fixed
$\IntegrationPoint\in\AmbientCarrier$,
we have
\begin{equation}\label{eq:riesz-pointwise-integral}
 (\AmbientProjection\HilbertFunction)(\IntegrationPoint)
 =\frac{1}{2\CircleConstant\ImaginaryUnit}
 \int_\SpectralContour
 \bigl[(\ResolventParameter\IdentityMap-\AmbientOperator)^{-1}
       \HilbertFunction\bigr](\IntegrationPoint)
 \,\IntegrationDifferential\ResolventParameter.
\end{equation}
The integration variable
$\ResolventParameter$
ranges over the boundary
$\SpectralContour$,
while the point
$\IntegrationPoint\in\AmbientCarrier$
is fixed.
Evaluation at
$\IntegrationPoint$
is a bounded linear functional on
$\ContinuousFunctions(\AmbientCarrier,\ComplexNumbers)$,
so it commutes with the line integral.
Equation \eqref{eq:riesz-pointwise-integral} also holds for
$\AmbientProjection_\SequenceIndex$
with
$\AmbientOperator_\SequenceIndex$
in place of
$\AmbientOperator$.

\textbf{Step 3. Convergence and dimensions of the ranges.}
\paraid{p-a982b8ced6ab}
By \Cref{lem:riesz-range-generalized-eigenspace},
we obtain
\begin{equation}\label{eq:ambient-projection-range}
 \range\AmbientProjection
 =\bigoplus_{\Eigenvalue\in\spec(\AmbientOperator)\cap\SpectralRegion}
   \GeneralizedEigenspace{\AmbientOperator}{\Eigenvalue}.
\end{equation}
For all sufficiently large
$\SequenceIndex$,
we also have
\begin{equation}\label{eq:approximating-projection-range}
 \range\AmbientProjection_\SequenceIndex
 =\bigoplus_{\Eigenvalue\in\spec(\AmbientOperator_\SequenceIndex)\cap\SpectralRegion}
   \GeneralizedEigenspace{\AmbientOperator_\SequenceIndex}{\Eigenvalue}.
\end{equation}
The sums are finite by \Cref{thm:compact-spectrum},
since
$0\notin\overline\SpectralRegion$.
By \Cref{lem:ambient-eigenspaces},
for every
$\Eigenvalue\in\spec(\AmbientOperator)\cap\SpectralRegion$,
\begin{equation}\label{eq:projection-range-eigenspace}
 \GeneralizedEigenspace{\AmbientOperator}{\Eigenvalue}
 =\ker(\AmbientOperator-\Eigenvalue\IdentityMap).
\end{equation}
For every sufficiently large
$\SequenceIndex$
and every
$\Eigenvalue\in\spec(\AmbientOperator_\SequenceIndex)\cap\SpectralRegion$,
the same proposition implies
\begin{equation}\label{eq:approximating-range-eigenspace}
 \GeneralizedEigenspace{\AmbientOperator_\SequenceIndex}{\Eigenvalue}
 =\ker(\AmbientOperator_\SequenceIndex-\Eigenvalue\IdentityMap).
\end{equation}
We can now apply \Cref{thm:riesz-convergence}
on the fixed Banach space
$\ContinuousFunctions(\AmbientCarrier,\ComplexNumbers)$.
The hypotheses on operators in the statement of
\Cref{thm:riesz-convergence}
follow from \eqref{eq:ambient-strong-convergence}
and the collective compactness of
$\{\AmbientOperator_\SequenceIndex-\AmbientOperator\}_{\SequenceIndex\in\NonnegativeIntegers}$.
Equations \eqref{eq:common-spectral-contour},
\eqref{eq:spectral-region-selection-1},
\eqref{eq:spectral-region-selection-2} and \eqref{eq:spectral-region-selection-3}
verify its hypotheses on
$\SpectralContour$.
For all sufficiently large
$\SequenceIndex$,
we obtain
\begin{equation}\label{eq:spectral-projection-convergence-1}
 \dim\range \AmbientProjection _\SequenceIndex =\dim\range \AmbientProjection .
\end{equation}
For every
$\HilbertFunction\in\ContinuousFunctions(\AmbientCarrier,\ComplexNumbers)$,
we also obtain
\begin{equation}\label{eq:spectral-projection-convergence-2}
 \norm{\AmbientProjection _\SequenceIndex \HilbertFunction -\AmbientProjection \HilbertFunction }_\infty\to0.
\end{equation}
The dimensions are finite because
$0\notin\overline\SpectralRegion $.
The convergence in \eqref{eq:spectral-projection-convergence-2} means
uniform convergence on
$\AmbientCarrier$
for each fixed function
$\HilbertFunction\in\ContinuousFunctions(\AmbientCarrier,\ComplexNumbers)$.

\paraid{p-60dd00c4dbee}
We now compare the projections on
$\ContinuousFunctions(\AmbientCarrier,\ComplexNumbers)$
with the spectral subspaces of the distance operators.
For
$\ComparisonCarrier=\BaseCarrier$
or
$\ComparisonCarrier=\BaseCarrier_\SequenceIndex$
with
$\SequenceIndex$
sufficiently large,
write
$\AmbientProjection_\BaseCarrier=\AmbientProjection$
and
$\AmbientProjection_{\BaseCarrier_\SequenceIndex}=\AmbientProjection_\SequenceIndex$.
By \eqref{eq:spectral-region-selection-1} and
\eqref{eq:spectral-region-selection-2},
$\SpectralRegion$
contains exactly those $\Eigenvalue\in\spec(\AmbientOperator)$ with
$|\Eigenvalue|>\SpectralCutoff$,
and exactly those $\Eigenvalue\in\spec(\AmbientOperator_\SequenceIndex)$ with
$|\Eigenvalue|>\SpectralCutoff$.
\Cref{lem:projection-restriction} induces the complex-linear isomorphism
\[
 \RestrictionOperator_\ComparisonCarrier\colon
 \range\AmbientProjection_\ComparisonCarrier
 \to
 \SpectralSubspace{\ComparisonCarrier}{\SpectralCutoff}
 +\ImaginaryUnit\SpectralSubspace{\ComparisonCarrier}{\SpectralCutoff}.
\]
Here the space on the right is the complexification of the real space
$\SpectralSubspace{\ComparisonCarrier}{\SpectralCutoff}$.
Its complex dimension equals
$\dim_\RealNumbers\SpectralSubspace{\ComparisonCarrier}{\SpectralCutoff}$.
The dimension equality in \eqref{eq:spectral-projection-convergence-1}
therefore yields
\begin{equation}\label{eq:spectral-dimension-identification}
 \begin{aligned}
 \dim_\RealNumbers\SpectralSubspace{\BaseCarrier_\SequenceIndex}{\SpectralCutoff}
 &=\dim_\ComplexNumbers\range\AmbientProjection_\SequenceIndex
 =\dim_\ComplexNumbers\range\AmbientProjection\\
 &=\dim_\RealNumbers\SpectralSubspace{\BaseCarrier}{\SpectralCutoff}
 =\SpectralRank.
 \end{aligned}
\end{equation}
\paraid{p-7a7a27ecec95}
Finally, applying \Cref{lem:projection-restriction} to the limit carrier
and to each sufficiently late carrier $\BaseCarrier_\SequenceIndex$ shows that
$\AmbientProjection$ and $\AmbientProjection_\SequenceIndex$ preserve real
functions.
\end{proof}

\begin{theorem}\label{lem:spectral-continuity}
\paraid{p-56449b80711e}
Let
$\BaseCarrier _\SequenceIndex ,\BaseCarrier $
be compact subsets of a compact metric space
$\MetricSpace{\AmbientCarrier }{\AmbientMetric }$,
with restricted metrics
\[
 \MetricSymbol_\SequenceIndex=\AmbientMetric|_{\BaseCarrier_\SequenceIndex\times\BaseCarrier_\SequenceIndex},
\]
and
\[
 \MetricSymbol=\AmbientMetric|_{\BaseCarrier\times\BaseCarrier},
\]
such that
\[
 \HausdorffDistance{\AmbientMetric }(\BaseCarrier _\SequenceIndex ,\BaseCarrier )\to0,
\]
and let $\iota_\SequenceIndex\colon\BaseCarrier_\SequenceIndex\hookrightarrow\AmbientCarrier$
and $\iota\colon\BaseCarrier\hookrightarrow\AmbientCarrier$ be the inclusions.
Define probabilities on $\AmbientCarrier$ by
$\bar\mu_\SequenceIndex=(\iota_\SequenceIndex)_*\SelectedMeasure{\BaseCarrier_\SequenceIndex}{\MetricSymbol_\SequenceIndex}$
and $\bar\mu=\iota_*\SelectedMeasure{\BaseCarrier}{\MetricSymbol}$.
Fix
$\SpectralCutoff >0$
with
$\SpectralCutoff ,-\SpectralCutoff \notin\spec(\DistanceOperator{\BaseCarrier}{\MetricSymbol})$.
Put
$\SpectralRank=\dim\SpectralSubspace{\BaseCarrier}{\SpectralCutoff}$.
There exists $\SequenceIndex_0\in\NonnegativeIntegers$ such that, for every
$\SequenceIndex\geq\SequenceIndex_0$,
\begin{equation}\label{eq:stable-spectral-endpoints}
 \SpectralCutoff,-\SpectralCutoff
 \notin\spec(\DistanceOperator{\BaseCarrier_\SequenceIndex}{\MetricSymbol_\SequenceIndex}),
\end{equation}
and
\begin{equation}\label{eq:stable-spectral-rank}
 \dim\SpectralSubspace{\BaseCarrier_\SequenceIndex}{\SpectralCutoff}
 =\SpectralRank.
\end{equation}
If
$\SpectralRank>0$,
fix any real orthonormal basis
$\{\Eigenfunction_\BasisIndex\}_{0\leq\BasisIndex<\SpectralRank}$
of
$\SpectralSubspace{\BaseCarrier}{\SpectralCutoff}$.
For the fixed basis
$\{\Eigenfunction_\BasisIndex\}_{0\leq\BasisIndex<\SpectralRank}$
and the sequence
$\{(\BaseCarrier_\SequenceIndex,\bar\mu_\SequenceIndex)\}_{\SequenceIndex\in\NonnegativeIntegers}$,
increase
$\SequenceIndex_0$
if necessary.
Then there exist families of functions
\[
 \{\Eigenfunction_{\SequenceIndex,\BasisIndex}
   \colon\BaseCarrier_\SequenceIndex\to\RealNumbers\}_{\SequenceIndex\geq\SequenceIndex_0,\,0\leq\BasisIndex<\SpectralRank},
\]
\[
 \{\BasisExtension_{\SequenceIndex,\BasisIndex}
   \colon\AmbientCarrier\to\RealNumbers\}_{\SequenceIndex\geq\SequenceIndex_0,\,0\leq\BasisIndex<\SpectralRank},
\]
and
\[
 \{\BasisExtension_\BasisIndex
   \colon\AmbientCarrier\to\RealNumbers\}_{0\leq\BasisIndex<\SpectralRank}
\]
with the following properties.
\begin{enumerate}[label=\textup{(B\arabic*)},ref=\textup{(B\arabic*)}]
\item\label{item:spectral-basis-orthonormality}
For every
$\SequenceIndex\geq\SequenceIndex_0$,
the family
$\{\Eigenfunction_{\SequenceIndex,\BasisIndex}\}_{0\leq\BasisIndex<\SpectralRank}$
is a real orthonormal basis of
$\SpectralSubspace{\BaseCarrier_\SequenceIndex}{\SpectralCutoff}$.
\item\label{item:spectral-basis-extensions}
For every
$\SequenceIndex\geq\SequenceIndex_0$
and
$0\leq\BasisIndex<\SpectralRank$,
the functions
$\BasisExtension_{\SequenceIndex,\BasisIndex},\BasisExtension_\BasisIndex
 \colon\AmbientCarrier\to\RealNumbers$
are continuous and satisfy
\[
 \BasisExtension_{\SequenceIndex,\BasisIndex}|_{\BaseCarrier_\SequenceIndex}
 =\Eigenfunction_{\SequenceIndex,\BasisIndex},
\]
and
\[
 \BasisExtension_\BasisIndex|_\BaseCarrier=\Eigenfunction_\BasisIndex.
\]
\item\label{item:spectral-basis-uniform-convergence}
The extensions also satisfy
\begin{equation}\label{eq:aligned-basis-extensions}
 \max_{0\leq\BasisIndex<\SpectralRank}
 \norm{\BasisExtension_{\SequenceIndex,\BasisIndex}-\BasisExtension_\BasisIndex}_\infty
 \longrightarrow0.
\end{equation}
\item\label{item:spectral-basis-correspondence-convergence}
For every sequence of correspondences
$\{\Correspondence_\SequenceIndex\subset\BaseCarrier_\SequenceIndex\times\BaseCarrier\}_{\SequenceIndex\geq\SequenceIndex_0}$
satisfying
\[
 \sup_{(\BasePoint_\SequenceIndex,\BasePoint)\in\Correspondence_\SequenceIndex}
 \AmbientMetric(\BasePoint_\SequenceIndex,\BasePoint)\longrightarrow0,
\]
the bases
$\{\Eigenfunction_{\SequenceIndex,\BasisIndex}\}_{0\leq\BasisIndex<\SpectralRank}$
and
$\{\Eigenfunction_\BasisIndex\}_{0\leq\BasisIndex<\SpectralRank}$
satisfy
\begin{equation}\label{eq:aligned-bases}
 \max_{0\leq\BasisIndex<\SpectralRank}
 \sup_{(\BasePoint_\SequenceIndex,\BasePoint)\in\Correspondence_\SequenceIndex}
 |\Eigenfunction_{\SequenceIndex,\BasisIndex}(\BasePoint_\SequenceIndex)
       -\Eigenfunction_\BasisIndex(\BasePoint)|\longrightarrow0.
\end{equation}
\end{enumerate}
\end{theorem}

\begin{proof}
\paraid{p-dc3fbce63b61}
By \Cref{thm:invariant-measures}\ref{item:invariant-measures-continuity} applied to the inclusions
$\iota_\SequenceIndex$
and
$\iota$,
we have
$\bar\mu_\SequenceIndex\to\bar\mu$
weakly.
By \Cref{lem:spectral-projections},
we obtain \eqref{eq:stable-spectral-endpoints}
and \eqref{eq:stable-spectral-rank}.
Let
$\AmbientProjection_\SequenceIndex,\AmbientProjection$
be the projections constructed in that proposition.
If
$\SpectralRank=0$,
the stable spectral subspaces are zero and no basis or extension data are required.
Assume that
$\SpectralRank>0$.
To construct the real orthonormal bases
$\{\Eigenfunction_{\SequenceIndex,\BasisIndex}\}_{0\leq\BasisIndex<\SpectralRank}$
of
$\SpectralSubspace{\BaseCarrier_\SequenceIndex}{\SpectralCutoff}$
in \ref{item:spectral-basis-orthonormality},
we first project continuous extensions of the prescribed basis
$\{\Eigenfunction_\BasisIndex\}_{0\leq\BasisIndex<\SpectralRank}$.
The resulting functions need not be orthonormal for the varying measures
$\SelectedMeasure{\BaseCarrier_\SequenceIndex}{\MetricSymbol_\SequenceIndex}$.
We use the Gram matrix to normalize these functions without changing their limit.
We then verify \ref{item:spectral-basis-extensions}
and \ref{item:spectral-basis-uniform-convergence}
for the normalized extensions and deduce
\ref{item:spectral-basis-correspondence-convergence}
from uniform continuity.

\ProofStep{Projected extensions and convergence of the Gram matrices.}\label{step:spectral-basis-projection}
\paraid{p-b8f758726161}
For every integer
$0\leq\BasisIndex<\SpectralRank$,
use the Tietze--Urysohn theorem (\Cref{thm:tietze}) to choose a continuous real extension
$\widetilde\Eigenfunction_\BasisIndex\colon\AmbientCarrier\to\RealNumbers$
of the prescribed function
$\Eigenfunction_\BasisIndex\colon\BaseCarrier\to\RealNumbers$.
For every sufficiently large
$\SequenceIndex\in\NonnegativeIntegers$
and every integer
$0\leq\BasisIndex<\SpectralRank$,
put
\[
 \ProjectedBasisFunction _{\SequenceIndex ,\BasisIndex }=\AmbientProjection _\SequenceIndex \widetilde \Eigenfunction _\BasisIndex .
\]
For every integer
$0\leq\BasisIndex<\SpectralRank$,
put
\[
 \ProjectedBasisFunction _\BasisIndex =\AmbientProjection \widetilde \Eigenfunction _\BasisIndex .
\]
Applying \eqref{eq:spectral-projection-convergence-2} to each fixed function
$\widetilde\Eigenfunction_\BasisIndex$,
we obtain
$\norm{\ProjectedBasisFunction _{\SequenceIndex ,\BasisIndex }-\ProjectedBasisFunction _\BasisIndex }_\infty\to0$
for every
$0\leq\BasisIndex<\SpectralRank$.
Moreover,
\eqref{eq:projection-restriction} implies
\[
 \RestrictionOperator_\BaseCarrier\ProjectedBasisFunction_\BasisIndex
 =\SpectralProjection{\BaseCarrier}{\SpectralCutoff}
   \RestrictionOperator_\BaseCarrier\widetilde\Eigenfunction_\BasisIndex
 =\SpectralProjection{\BaseCarrier}{\SpectralCutoff}\Eigenfunction_\BasisIndex
 =\Eigenfunction_\BasisIndex.
\]
Thus
$\ProjectedBasisFunction_\BasisIndex|_\BaseCarrier=\Eigenfunction_\BasisIndex$.
Equation \eqref{eq:projection-restriction} on
$\BaseCarrier_\SequenceIndex$
shows that
\[
 \ProjectedBasisFunction_{\SequenceIndex,\BasisIndex}|_{\BaseCarrier_\SequenceIndex}
 =\SpectralProjection{\BaseCarrier_\SequenceIndex}{\SpectralCutoff}
   (\widetilde\Eigenfunction_\BasisIndex|_{\BaseCarrier_\SequenceIndex})
 \in\SpectralSubspace{\BaseCarrier_\SequenceIndex}{\SpectralCutoff}.
\]
Using
$\IntegrationPoint\in\AmbientCarrier$
as the integration variable,
define the real Gram matrix
\[
 \GramMatrix _\SequenceIndex =\left(\int_\AmbientCarrier  \ProjectedBasisFunction _{\SequenceIndex ,\BasisIndex }(\IntegrationPoint)\ProjectedBasisFunction _{\SequenceIndex ,\SecondBasisIndex }(\IntegrationPoint)\,\IntegrationDifferential \bar\mu_\SequenceIndex(\IntegrationPoint)
                                     \right)_{0\leq \BasisIndex ,\SecondBasisIndex<\SpectralRank }.
\]
For every pair of integers
$0\leq\BasisIndex,\SecondBasisIndex<\SpectralRank$,
the products
$\ProjectedBasisFunction_{\SequenceIndex,\BasisIndex}
 \ProjectedBasisFunction_{\SequenceIndex,\SecondBasisIndex}$
converge uniformly on
$\AmbientCarrier$
to
$\ProjectedBasisFunction_\BasisIndex\ProjectedBasisFunction_\SecondBasisIndex$.
Since
$\ProjectedBasisFunction_\BasisIndex|_\BaseCarrier=\Eigenfunction_\BasisIndex$
and the prescribed basis is orthonormal,
\[
 \int_\AmbientCarrier
 \ProjectedBasisFunction_\BasisIndex(\IntegrationPoint)\ProjectedBasisFunction_\SecondBasisIndex(\IntegrationPoint)
 \,\IntegrationDifferential\bar\mu(\IntegrationPoint)
 =\KroneckerSymbol_{\BasisIndex\SecondBasisIndex}.
\]
For every sufficiently large
$\SequenceIndex$
and every pair of integers
$0\leq\BasisIndex,\SecondBasisIndex<\SpectralRank$,
define the uniform product error by
\begin{equation}\label{eq:gram-uniform-product-error}
 a_{\SequenceIndex,\BasisIndex,\SecondBasisIndex}
 =\norm{\ProjectedBasisFunction_{\SequenceIndex,\BasisIndex}
       \ProjectedBasisFunction_{\SequenceIndex,\SecondBasisIndex}
       -\ProjectedBasisFunction_\BasisIndex\ProjectedBasisFunction_\SecondBasisIndex}_\infty.
\end{equation}
For the same indices,
define the integration error for the fixed limiting product by
\begin{equation}\label{eq:gram-weak-integration-error}
 b_{\SequenceIndex,\BasisIndex,\SecondBasisIndex}
 =\left|\int_\AmbientCarrier
       \ProjectedBasisFunction_\BasisIndex(\IntegrationPoint)\ProjectedBasisFunction_\SecondBasisIndex(\IntegrationPoint)
       \,\IntegrationDifferential\bar\mu_\SequenceIndex(\IntegrationPoint)
       -\int_\AmbientCarrier
       \ProjectedBasisFunction_\BasisIndex(\IntegrationPoint)\ProjectedBasisFunction_\SecondBasisIndex(\IntegrationPoint)
       \,\IntegrationDifferential\bar\mu(\IntegrationPoint)\right|.
\end{equation}
Since
$\bar\mu_\SequenceIndex(\AmbientCarrier)=1$,
\begin{equation}\label{eq:gram-entry-error}
 \left|(\GramMatrix_\SequenceIndex)_{\BasisIndex\SecondBasisIndex}
             -\KroneckerSymbol_{\BasisIndex\SecondBasisIndex}\right|
 \leq a_{\SequenceIndex,\BasisIndex,\SecondBasisIndex}
       +b_{\SequenceIndex,\BasisIndex,\SecondBasisIndex}.
\end{equation}
For every fixed
$0\leq\BasisIndex,\SecondBasisIndex<\SpectralRank$,
uniform convergence of the products implies
$a_{\SequenceIndex,\BasisIndex,\SecondBasisIndex}\to0$
in \eqref{eq:gram-uniform-product-error},
and weak convergence of
$\bar\mu_\SequenceIndex$
to
$\bar\mu$
implies
$b_{\SequenceIndex,\BasisIndex,\SecondBasisIndex}\to0$
in \eqref{eq:gram-weak-integration-error}.
By \eqref{eq:gram-entry-error},
\[
 (\GramMatrix _\SequenceIndex )_{\BasisIndex \SecondBasisIndex }\to\int_\AmbientCarrier  \ProjectedBasisFunction _\BasisIndex(\IntegrationPoint) \ProjectedBasisFunction _\SecondBasisIndex(\IntegrationPoint) \,\IntegrationDifferential \bar\mu(\IntegrationPoint)=\KroneckerSymbol _{\BasisIndex \SecondBasisIndex }.
\]
Thus
$\GramMatrix _\SequenceIndex \to \IdentityMatrix _\SpectralRank $,
so choose
$\SequenceIndex_0\in\NonnegativeIntegers$
such that for every
$\SequenceIndex\geq\SequenceIndex_0$
the matrix
$\GramMatrix_\SequenceIndex$
is positive definite.

\ProofStep{Normalization by inverse square roots.}\label{step:spectral-basis-normalization}
\paraid{p-dc0f659821c3}
We prove that the inverse square roots of the Gram matrices converge to the identity.
\Cref{thm:matrix-square-root} defines
$\GramMatrix_\SequenceIndex^{-1/2}$
by replacing each positive eigenvalue with its inverse square root in
an orthonormal eigenbasis.
For a real
$\SpectralRank\times\SpectralRank$
matrix
$\MatrixOperator$,
write its Euclidean operator norm as
\[
 \norm{\MatrixOperator}_\OperatorNorm
 =\sup_{\substack{\CoefficientVector\in\RealNumbers^\SpectralRank\\
                  \norm{\CoefficientVector}_2=1}}
       \norm{\MatrixOperator\CoefficientVector}_2.
\]
For every sufficiently large
$\SequenceIndex\in\NonnegativeIntegers$,
put
\[
 \MatrixError_\SequenceIndex
 =\norm{\GramMatrix_\SequenceIndex-\IdentityMatrix_\SpectralRank}_\OperatorNorm.
\]
Entrywise convergence  implies
$\MatrixError_\SequenceIndex\to0$.
Increase
$\SequenceIndex_0$
if necessary so that for every
$\SequenceIndex\geq\SequenceIndex_0$,
\[
 0\leq\MatrixError_\SequenceIndex<1.
\]
Fix
$\SequenceIndex\geq\SequenceIndex_0$.
Since
$\GramMatrix_\SequenceIndex$
is real symmetric,
\Cref{thm:matrix-square-root} implies that every
$\Eigenvalue\in\spec(\GramMatrix_\SequenceIndex)$
is real and admits an eigenvector
$\CoefficientVector\in\RealNumbers^\SpectralRank$
with
$\norm{\CoefficientVector}_2=1$.
For such
$\Eigenvalue$
and
$\CoefficientVector$,
we have
\[
 |\Eigenvalue-1|
 =\norm{(\GramMatrix_\SequenceIndex-\IdentityMatrix_\SpectralRank)
         \CoefficientVector}_2
 \leq\norm{\GramMatrix_\SequenceIndex-\IdentityMatrix_\SpectralRank}_\OperatorNorm
      \norm{\CoefficientVector}_2
 =\MatrixError_\SequenceIndex.
\]
Since
$\Eigenvalue\in\RealNumbers$
and
$\MatrixError_\SequenceIndex<1$,
it follows that
\[
 \spec(\GramMatrix_\SequenceIndex)
 \subset[1-\MatrixError_\SequenceIndex,1+\MatrixError_\SequenceIndex]
 \subset(0,\infty).
\]
For these indices
$\SequenceIndex$,
\Cref{thm:matrix-square-root} implies
\[
 \begin{aligned}
 \norm{\GramMatrix_\SequenceIndex^{-1/2}-\IdentityMatrix_\SpectralRank}_\OperatorNorm
 &=\max_{\Eigenvalue\in\spec(\GramMatrix_\SequenceIndex)}
       |\Eigenvalue^{-1/2}-1|
 \leq\max_{1-\MatrixError_\SequenceIndex\leq\SecondParameter\leq1+\MatrixError_\SequenceIndex}
       |\SecondParameter^{-1/2}-1|\\
 &=\max_{1-\MatrixError_\SequenceIndex\leq\SecondParameter\leq1+\MatrixError_\SequenceIndex}
       \frac{|\SecondParameter-1|}
       {\sqrt{\SecondParameter}(1+\sqrt{\SecondParameter})}\\
 &\leq
 \frac{\MatrixError_\SequenceIndex}
 {\sqrt{1-\MatrixError_\SequenceIndex}
  (1+\sqrt{1-\MatrixError_\SequenceIndex})}
 \longrightarrow0.
 \end{aligned}
\]
Thus
\[
 \GramMatrix_\SequenceIndex^{-1/2}
 \xrightarrow[\SequenceIndex\to\infty]{\norm{\cdot}_\OperatorNorm}
 \IdentityMatrix_\SpectralRank.
\]
For every
$\SequenceIndex\geq\SequenceIndex_0$
and
$\IntegrationPoint\in\AmbientCarrier$,
we define the row map
\[
 \ProjectedBasisRow_\SequenceIndex\colon\AmbientCarrier\to\RealNumbers^{1\times\SpectralRank}
\]
by
\[
 \ProjectedBasisRow_\SequenceIndex(\IntegrationPoint)
 =(\ProjectedBasisFunction_{\SequenceIndex,0}(\IntegrationPoint),\ldots,
   \ProjectedBasisFunction_{\SequenceIndex,\SpectralRank-1}(\IntegrationPoint)).
\]
We multiply this row vector by the matrix
$\GramMatrix_\SequenceIndex^{-1/2}$
on the right and define
\begin{equation}\label{eq:normalized-basis-row}
 \bigl(\widetilde\Eigenfunction_{\SequenceIndex,0}(\IntegrationPoint),\ldots,
       \widetilde\Eigenfunction_{\SequenceIndex,\SpectralRank-1}(\IntegrationPoint)\bigr)
 =\ProjectedBasisRow_\SequenceIndex(\IntegrationPoint)\GramMatrix_\SequenceIndex^{-1/2}.
\end{equation}
For every integer
$0\leq\BasisIndex<\SpectralRank$,
the components in \eqref{eq:normalized-basis-row} satisfy
\[
 \widetilde\Eigenfunction_{\SequenceIndex,\BasisIndex}(\IntegrationPoint)
 =\sum_{\SecondBasisIndex=0}^{\SpectralRank-1}
 \ProjectedBasisFunction_{\SequenceIndex,\SecondBasisIndex}(\IntegrationPoint)
 (\GramMatrix_\SequenceIndex^{-1/2})_{\SecondBasisIndex\BasisIndex},
\]
and
\[
 \Eigenfunction_{\SequenceIndex,\BasisIndex}
 =\widetilde\Eigenfunction_{\SequenceIndex,\BasisIndex}|_{\BaseCarrier_\SequenceIndex}.
\]
Since
$\GramMatrix_\SequenceIndex^{-1/2}$
is symmetric,
the Gram matrix of these restrictions is
\begin{equation}\label{eq:normalized-basis-gram}
 \begin{aligned}
 &\int_\AmbientCarrier
   \bigl(\ProjectedBasisRow_\SequenceIndex(\IntegrationPoint)
         \GramMatrix_\SequenceIndex^{-1/2}\bigr)^{\TransposeSymbol}
   \bigl(\ProjectedBasisRow_\SequenceIndex(\IntegrationPoint)
         \GramMatrix_\SequenceIndex^{-1/2}\bigr)
   \,\IntegrationDifferential\bar\mu_\SequenceIndex(\IntegrationPoint)\\
 &=(\GramMatrix_\SequenceIndex^{-1/2})^{\TransposeSymbol}
   \GramMatrix_\SequenceIndex\GramMatrix_\SequenceIndex^{-1/2}\\
 &=\GramMatrix_\SequenceIndex^{-1/2}\GramMatrix_\SequenceIndex
   \GramMatrix_\SequenceIndex^{-1/2}=\IdentityMatrix_\SpectralRank.
 \end{aligned}
\end{equation}
By \eqref{eq:projection-restriction},
the functions
$\Eigenfunction_{\SequenceIndex,\BasisIndex}$
belong to
$\SpectralSubspace{\BaseCarrier_\SequenceIndex}{\SpectralCutoff}$.
Equation \eqref{eq:normalized-basis-gram} proves orthonormality,
and their number equals
$\dim\SpectralSubspace{\BaseCarrier_\SequenceIndex}{\SpectralCutoff}$,
so they form an orthonormal basis of
$\SpectralSubspace{\BaseCarrier_\SequenceIndex}{\SpectralCutoff}$.
This proves \ref{item:spectral-basis-orthonormality}.

\ProofStep{Continuous extensions and convergence.}\label{step:spectral-basis-extensions}
\paraid{p-9a01e68ec997}
For every
$\SequenceIndex\geq\SequenceIndex_0$
and every integer
$0\leq\BasisIndex<\SpectralRank$,
put
$\BasisExtension_{\SequenceIndex,\BasisIndex}
 =\widetilde\Eigenfunction_{\SequenceIndex,\BasisIndex}$
and
$\BasisExtension_\BasisIndex=\ProjectedBasisFunction_\BasisIndex$.
The projections preserve real functions,
so the construction in Step~\ref{step:spectral-basis-projection} and the finite linear combinations in
\eqref{eq:normalized-basis-row} show that these functions are real and continuous.
Their restriction identities follow from Step~\ref{step:spectral-basis-projection} and the definition of
$\Eigenfunction_{\SequenceIndex,\BasisIndex}$
in Step~\ref{step:spectral-basis-normalization}.
This proves \ref{item:spectral-basis-extensions}.
We next prove \ref{item:spectral-basis-uniform-convergence}.
The uniform convergence of the projected basis functions in Step~\ref{step:spectral-basis-projection},
the convergence
$\GramMatrix_\SequenceIndex^{-1/2}\to\IdentityMatrix_\SpectralRank$,
and \eqref{eq:normalized-basis-row} imply that,
for every integer
$0\leq\BasisIndex<\SpectralRank$,
\[
 \norm{\widetilde\Eigenfunction_{\SequenceIndex,\BasisIndex}
       -\ProjectedBasisFunction_\BasisIndex}_\infty\longrightarrow0.
\]
Since there are finitely many basis indices,
it follows that
\[
 \max_{0\leq\BasisIndex<\SpectralRank}
 \norm{\widetilde\Eigenfunction_{\SequenceIndex,\BasisIndex}
       -\ProjectedBasisFunction_\BasisIndex}_\infty\longrightarrow0.
\]
By the definitions of the extensions,
\[
 \max_{0\leq\BasisIndex<\SpectralRank}
 \norm{\BasisExtension_{\SequenceIndex,\BasisIndex}
       -\BasisExtension_\BasisIndex}_\infty
 =\max_{0\leq\BasisIndex<\SpectralRank}
 \norm{\widetilde\Eigenfunction_{\SequenceIndex,\BasisIndex}
       -\ProjectedBasisFunction_\BasisIndex}_\infty
 \longrightarrow0,
\]
which proves \ref{item:spectral-basis-uniform-convergence}.
Finally,
let
$\Correspondence_\SequenceIndex\subset\BaseCarrier_\SequenceIndex\times\BaseCarrier$
be a sequence of correspondences satisfying the hypothesis in
\ref{item:spectral-basis-correspondence-convergence}.
For every
$\SequenceIndex\geq\SequenceIndex_0$,
put
\[
 \CorrespondenceDisplacement_\SequenceIndex
 =\sup_{(\BasePoint_\SequenceIndex,\BasePoint)\in\Correspondence_\SequenceIndex}
       \AmbientMetric(\BasePoint_\SequenceIndex,\BasePoint).
\]
For every integer
$0\leq\BasisIndex<\SpectralRank$,
the restriction identities in \ref{item:spectral-basis-extensions} imply
\begin{equation}\label{eq:basis-alignment-error}
 \sup_{(\BasePoint_\SequenceIndex,\BasePoint)\in\Correspondence_\SequenceIndex}|\Eigenfunction _{\SequenceIndex ,\BasisIndex }(\BasePoint _\SequenceIndex )-\Eigenfunction _\BasisIndex (\BasePoint )|
 \leq\norm{\widetilde \Eigenfunction _{\SequenceIndex ,\BasisIndex }-\ProjectedBasisFunction _\BasisIndex }_\infty
       +\sup_{\substack{\IntegrationPoint ,\IntegrationPoint _0\in \AmbientCarrier \\\AmbientMetric (\IntegrationPoint ,\IntegrationPoint _0)\leq\CorrespondenceDisplacement _\SequenceIndex }}
                  |\ProjectedBasisFunction _\BasisIndex (\IntegrationPoint )-\ProjectedBasisFunction _\BasisIndex (\IntegrationPoint _0)|.
\end{equation}
The maximum over the basis indices of the norm term in
\eqref{eq:basis-alignment-error} tends to
$0$
by
\ref{item:spectral-basis-uniform-convergence}
and the definitions of the extensions.
For each
$0\leq\BasisIndex<\SpectralRank$,
uniform continuity of
$\ProjectedBasisFunction_\BasisIndex$
on the compact space $\AmbientCarrier$ and
$\CorrespondenceDisplacement_\SequenceIndex\to0$
imply the following convergence,
where the maximum tends to
$0$
because there are finitely many basis indices.
\[
 \max_{0\leq\BasisIndex<\SpectralRank}
 \sup_{\substack{\IntegrationPoint,\IntegrationPoint_0\in\AmbientCarrier\\
                  \AmbientMetric(\IntegrationPoint,\IntegrationPoint_0)
                  \leq\CorrespondenceDisplacement_\SequenceIndex}}
 \left|\ProjectedBasisFunction_\BasisIndex(\IntegrationPoint)
       -\ProjectedBasisFunction_\BasisIndex(\IntegrationPoint_0)\right|
 \longrightarrow0.
\]
Taking the maximum over the basis indices in
\eqref{eq:basis-alignment-error} now yields
\[
 \max_{0\leq\BasisIndex<\SpectralRank}
 \sup_{(\BasePoint_\SequenceIndex,\BasePoint)\in\Correspondence_\SequenceIndex}
 \left|\Eigenfunction_{\SequenceIndex,\BasisIndex}(\BasePoint_\SequenceIndex)
       -\Eigenfunction_\BasisIndex(\BasePoint)\right|
 \longrightarrow0,
\]
which proves \ref{item:spectral-basis-correspondence-convergence}.
\end{proof}
